%
% 41-definition.tex -- Definition
%
% (c) 2026 Prof Dr Andreas Müller
%
\subsection{Definition
\label{buch:analytisch:ganz:subsection:definition}}
Während meromorphe Funktionen nur in isolierten Punkten nicht definiert
sind, sind ganze Funktionen gemäss der folgenden Definition überall
definiert.

\begin{definition}[Ganze Funktion]
Eine holomorphe Funktion, die auf ganz $\mathbb{C}$ definiert ist,
heisst \emph{ganz}.
\index{ganze Funktion}%
\end{definition}

\begin{beispiel}
Jedes Polynom $a(z) = a_0 + a_1z + \ldots + a_{n-1}z^{n-1} + a_nz^n$ 
mit $n>0$ ist eine ganze Funktion.
\end{beispiel}

\begin{beispiel}
Die Exponentialfunktion
\[
e^z
=
1
+
z
+
\frac{z^2}{2}
+
\frac{z^3}{3!}
+
\frac{z^4}{4!}
+
\ldots
=
\sum_{k=0}^\infty \frac{z^k}{k!}
\]
ist eine ganze Funktion, denn ihr Konvergenzradius ist gemäss
Beispiel~\ref{buch:komplex:funktionen:bsp:exponentialfunktion}
unendlich gross.
Die aus $e^z$ abgeleiteten Funktionen $\sin z$ und $\cos z$ sind
ebenfalls ganz.
\end{beispiel}

